\documentclass[10pt,a4paper]{amsart}
\usepackage[margin=19mm]{geometry}
\usepackage{amsmath,amssymb,mathtools}
\usepackage[scr=rsfs,cal=euler]{mathalfa}
\usepackage{xfrac}
\usepackage[unicode,hidelinks,bookmarks=false]{hyperref}
\newcommand{\ie}{i.e.\ }
\newcommand{\cf}{cf.\ }
\newcommand{\resp}{respectively\ }
\newcommand{\Subsection}{\subsection}
\providecommand{\nobreakdash}{\nobreak\hbox{-}}
\setlength{\emergencystretch}{2em}
\multlinegap0pt
\usepackage{bm}
\usepackage{bbm}
\usepackage{dsfont}
\usepackage{xspace}
\usepackage{cancel}
\usepackage{mathtools}
\usepackage{xcolor}
\usepackage{amsthm}
\makeatletter
\@ifundefined{Theorem}{\newtheorem{Theorem}{Theorem}}{}
\@ifundefined{Lemma}{\newtheorem{Lemma}[Theorem]{Lemma}}{}
\makeatother
\title[Stability of input-output maps and their minimal realization]{Stability of input-output maps and their minimal realizations in state-linear, state-affine, LPV, and linear switched systems}
\author{Mih\'aly Petreczky, Juan-Pablo Ortega, Florian Rossmannek, B\'alint Dar\'oczy}
\date{Source: arXiv:2607.03849v1 (CC BY 4.0)}
\begin{document}
\maketitle

\noindent\emph{Source and licence.} Stability of input-output maps and their minimal realizations in state-linear, state-affine, LPV, and linear switched systems. Version arXiv:2607.03849v1 (CC BY 4.0), distributed under the Creative Commons Attribution 4.0 International licence, as recorded in the arXiv licence metadata for this exact version and shown on the abstract page https://arxiv.org/abs/2607.03849v1. Full rights notice: licenses/arxiv-2607.03849-CC-BY-4.0.txt.\par\smallskip

\noindent\textit{Editorial scope.} Counted unit: the Lemma whose source label is \texttt{lem:lpv:IRR-UD} in the article (source0.tex lines 1732--1744, source label \texttt{lem:lpv:IRR-UD}), delivered together with the article's own complete original proof of it (source0.tex lines 1746--1833, one \texttt{proof} environment of 88 lines). No mathematics is written, repaired or re-derived: statement and proof are the article's own text line for line. Required context retained uncounted: none. Every definition, display and label of those lines is retained unchanged. Omitted from this package and not redistributed: 1--1731, 1745--1745, 1834--2067. Nothing inside a retained excerpt is omitted or altered: every retained body is delivered exactly as the article's lines are written, so no omission receipt of any kind accompanies this package. Citations: the retained excerpts carry no citation command at all, so no bibliography entry of the article is transcribed and no reference is redistributed. Reference adaptation: none was needed and none was made; every reference inside the delivered unit resolves inside it, so no unresolved reference or citation is produced. Printed slips kept literal: no printed word, symbol, hypothesis or number of the counted statement or of its proof was altered. Layout and class: the article's own class is \texttt{ieeecolor} with packages including amsmath, amssymb, mathtools, enumerate, graphicx, textcomp, algorithm, algorithmicx, cite, mathrsfs, lcsys, subcaption, float, bm; the wrapper is the lane's pinned \texttt{amsart} editorial wrapper at 10pt on A4, which restates the theorem-like environments the retained text needs and adds the article's own macro definitions that the retained text needs (none), with extra wrapper packages (bm, bbm, dsfont, xspace, cancel, mathtools, xcolor). The delivered numbering is the unit's own. Assets: no retained span contains a figure environment, a graphics-inclusion command, a PDF-page inclusion, a table or a TikZ picture; no image file is copied into this package, the package declares no asset dependency and no image file is redistributed with it. Licence: version arXiv:2607.03849v1 of this article is distributed under the Creative Commons Attribution 4.0 International licence, as recorded in the arXiv licence metadata for this exact version and shown on the abstract page https://arxiv.org/abs/2607.03849v1. A full rights notice is delivered at licenses/arxiv-2607.03849-CC-BY-4.0.txt. Characters: every retained excerpt is pure ASCII, so the delivered engine, XeLaTeX with its default Latin Modern text font, reports no missing character.
  \begin{Lemma}
  \label{lem:lpv:IRR-UD}
  $H$ is IRR-UD (resp. IRR-UED \textcolor{black}{with rate $\beta$}) if and only if $\Psi_H$ is UD 
  (resp. UED \textcolor{black}{with rate $\beta$}). % resp. UAS).
  %%\begin{enumerate}
  %%\item $H$ is IRR-UD 
  %%
  %%$\iff$  $h \diamond \underline{p}$ is uniformly decaying for all $\underline{p} \in \mathbb{P}^{\infty}$, i.e., 
  %%\item $H$ is IRR-UED $\iff$  the impulse response $h \diamond \underline{p}$ is uniformly exponentially decaying for all $\underline{p} \in \mathbb{P}^{\infty}$, i.e., there exists $\beta \in (0,1)$ such that $\sup_{\underline{p} \in \mathbb{P}^{\infty}} \|(h \diamond \underline{p})(N+\tau,\tau)\|_2 < C\beta^N$ for some $C > 0$.
  %%\item $H$ is IRR-UAS $\iff$  the impulse response $h \diamond \underline{p}$ is uniformly absolutely summable for all $\underline{p} \in \mathbb{P}^{\infty}$, i.e., $\sup_{\underline{p} \in \mathbb{P}^{\infty}, N \in \mathbb{N}}\sum_{r=0}^{N-1} \|(h \diamond \underline{p})(N,r)\|_2$ is bounded and 
  %M$\lim_{n \rightarrow +\infty} \sup_{\underline{p} \in \mathbb{P}^{\infty}} \sum_{r=n}^{N+n-1} \|(h \diamond \underline{p})(N+n,r)\|_2=0$.
  %%\end{enumerate}
  \end{Lemma}

  \begin{proof}[Proof of Lemma \ref{lem:lpv:IRR-UD}]
   We show that 
      \begin{equation}
      \label{eq:pf:lpv1}
      \begin{split}
           %&
          %\sup_{\underline{p} \in \mathbb{P}^{\infty}} \|(h \diamond \underline{p})(K,l)\|_2=
          %\sup_{q_l,q_N \in Q, s \in Q^{N-l-1}} \|S_{q_{N},q_{l}}(s)\|_2 \\
          \sup_{\underline{p} \in \mathbb{P}^{\infty}}\!\! %\sum_{r=l}^{K} 
          \|(h \diamond \underline{p})(N+\tau,\tau)\|_2
          \!\! = \!\! 
          \sup_{s \in Q^{N+1}}  \|S(s)\|_2 
          %\!\! \sum_{r=l}^{K} \|S(s_{r:N})\|_2
      \end{split}
      \end{equation}
   for any \textcolor{black}{$N \ge 1,\tau \in \mathbb{N}$}.  %the following holds:
      %where $s_{r:N}=s[r+1] \cdots s[N]$ is the suffix of $s$ starting at position $r+1$.
      %In particular,  $K=\tau$ and $l=\tau$, 
      %$\sup_{\underline{p} \in \mathbb{P}^{\infty}} \|(h \diamond \underline{p})(N-1,\tau)\|_2=\sup_{s \in Q^{N}} 
      The statement of the lemma follows from \eqref{eq:pf:lpv1}
      \textcolor{black}{by noticing that 
       $\frac{1}{\sqrt{m}} \|S(qwq_1)\|_2  \le \sup_{l=1,\ldots,m} \|M_{q,l}(w)\|_2 \le \sqrt{n_p} \sup_{\sigma  \in Q} \|S(q w\sigma)\|_2$, $q,q_1 \in Q,w \in Q^{*}$.}
      %this particular case, 
      %while the statement for UAS follows from the general case of }.   
     % For the latter, \eqref{eq:pf:lpv1} applied to $l=0$ and $K=T-1$, $N=T+n-1$ implies
     % notice that with the notation of \eqref{eq:IRR-UAS}
     % \[ 
     %     b_n=\sup_{T \in \mathbb{N}} \sup_{s \in Q^{T+n+1}}
     %      \sum_{r=0}^{T-1} \|S(s_{r:T+n})\|_2
     % \]
      % In turn, notice that for any 
        %    %$q_0,q_1 \ldots, \in Q$, 
        %       $\bar{q}_i \in Q$, $i \in \mathbb{N}$,  and any $T$      
        %       by setting $q_i=\bar{q}_{T+n-i}$ for $i=1,\ldots, T+n-1$, 
        %       and $s=q_0 \cdots q_{T+n}$,
        %       we have
        %    \[
        %        \begin{split}
        %         \sum_{r=n+1}^{T+n-1} \|M_{q,j}(\bar{q}_{r}\cdots \bar{q}_{1})\|_2= 
        %         \sum_{r=0}^{T-1} \|M_{q_r,j}(s_{r+1:T+n-1})\|_2 
        %       \end{split}    
        %       \]
        %       By uniform convergence of 
        %       $\sum_{r=0}^{\infty} \|M_{q,j}(\bar{q}_{r}\cdots \bar{q}_{1})\|_2$
        %       it follows that $\epsilon > 0$, there exists $K_{\epsilon}$  (independent of $\{\bar{q}_i\}_{i=0}^{\infty},q,j$) such that for any $n > K_{\epsilon}$,
        %       \(
        %          \sum_{r=n+1}^{N+n-1} \|M_{q,j}(\bar{q}_{r}\cdots \bar{q}_{1})\|_2 < \epsilon/m
        %       \), and hence 
        %       \[ \sum_{r=0}^{T-1} \|S(s_{r:T+n})\|_2 \le \sum_{j=1}^{m}\sum_{r=0}^{T-1} \|M_{q_r,j}(s_{r+1:T+n-1})\|_2 <\epsilon
        %       \]
        %       for any $T$ and any
        %       choice of $s \in Q^{T+n+1}$, which implies that $\lim_{n \rightarrow +\infty} b_n=0$.    
        To show \eqref{eq:pf:lpv1}, for any choice of $\underline{q} \in Q^{\infty}$, and any $N$, choose $\underline{p} \in \mathbb{P}^{\infty}$ such that $\underline{p}[i]=e_{\underline{q}[i]}$ \textcolor{black}{for all $i \ge 1$}. Then 
        $(h \diamond \underline{p})(N+\tau,\tau)=S(\textcolor{black}{\underline{q}[\tau] \cdots \underline{q}[N+\tau]})$
        and hence the left-hand side of \eqref{eq:pf:lpv1} is greater than or equal to the right-hand side.
        Conversely, for any choice of $\underline{p} \in \mathbb{P}^{\infty}$
	  by using the fact that $\mathbb{P} \subseteq \Delta$\textcolor{red}{,}  for any $N$, % is the standard simplex, for any $N$, 
        \[ 
          \begin{split}
             & \|(h \diamond \underline{p})(N+\tau,\tau)\|_2  \le \sum_{s \in Q^{N+1}} \|S(s)\|_2 \underline{p}_s(N+\tau) \\
            & \le \sup_{s \in Q^{N+1}} \|S(s)\|_2 \sum_{s \in Q^{N+1}} \underline{p}_s(N+\tau) =
            \sup_{s \in Q^{N+1}} \|S(s)\|_2 \textcolor{black}{,}
          \end{split}
        \]
        where we used that  $\sum_{s \in Q^l} \underline{p}_s(t) = 1$ for any $t \ge l$.
        %and that $\underline{p}_{s_1}(t-|s_2|)\underline{p}_{s_2}(t)=\underline{p}_{s_1 s_2}(t)$.
  %%    Hence,
  %%    %\[
  %%    %  \begin{split}
  %%    %  &  \|h \diamond \underline{p}(N,\tau)\|_2  \le \sup_{s \in Q^{N}} 
  %%    %     \|S(s)\|_2 \sum_{s \in Q^{N}} p_s(N) =
  %%    %     \sup_{s \in Q^{N+1}} \|S(s)\|_2 
  %%    %     %\times  \\ & \times \prod_{i=1}^{N+k} p_{q_i}(i) = 
  %%    %    %\sup_{q_\tau,q_{N} \in Q,s \in Q^{N-\tau-1}} 
  %%    %     %\|S_{q_{N}, q_r}(s)\|_2 
  %%    %  \end{split}
  %%    %\]
  %%    \[
  %%      \begin{split}
  %%      & \sum_{r=l}^{K} \|h \diamond \underline{p}(N-1,r)\|_2 \le \sum_{s \in Q^{N}} 
  %%         \sum_{r=l}^{K}\|S(s_{r:N})\|_2 p_s(N) = \\
  %%         %\prod_{i=1}^{N} p_{q_i}(i) = \\
  %%        &  \le \sup_{s \in Q^{N}} \sum_{r=l}^{K} \|S(s_{r:N})\|_2 
  %%        \!\! \sum_{s \in Q^{N}}\!\! p_s(N) \!\!=\!\! 
  %%         \sup_{s \in Q^{N}} \sum_{r=l}^{K-1} \|S(s_{r:N})\|_2 
  %%      \end{split}
  %%    \]
  \end{proof}

\end{document}
